\section{Optimización y GRATE: de modelos grandes a modelos compactos}
\subsection{Pipeline de entrenamiento con dos agentes y una métrica de calidad}
En el caso de 00:24:31--00:31:15, ServiceNow primero usa el modelo grande Llama 3.5, con 405B de parámetros, para producir interacciones de alta calidad, y después usa los registros de tape como datos para hacer fine-tuning de un modelo Llama de apenas 8B de parámetros y así recuperar el desempeño. El resultado es que el modelo ajustado obtiene una puntuación GRATE equivalente a la del modelo grande original, pero logra un ahorro de más de 300x en el costo de operación.

\begin{knowledgebox}{Puntuación de calidad GRATE}
GRATE es un conjunto de atributos de calidad cuantificables (grounded, responsive, accurate, transparent, effective); cada dimensión puede medirse y optimizarse de manera independiente, lo que facilita usar el tape como señal de reward o de regresión durante el entrenamiento.
\end{knowledgebox}

\subsection{Compromiso entre costo y desempeño}
La gráfica de métricas de 00:30:25--00:31:38 muestra: el eje x representa el costo por millón de turnos de conversación, y el eje y representa la puntuación GRATE; los modelos grandes siguen en la esquina superior derecha, pero los modelos ajustados generados mediante tape pueden alcanzar una calidad cercana con apenas una milésima parte del costo.

\subsection{Resumen de la sección}
TapeAgents no solo facilita el desarrollo con ingeniería, sino que también nos permite convertir un agente «maestro» costoso en un «agente estudiante» de alta calidad y bajo costo; la curva de costo y desempeño puede visualizarse mediante GRATE.

